\titledquestion{Intégration multivariée}
Calculez l'intégrale double suivante :
$$\iint_D xy \, dy \, dx$$
où $D$ est le domaine défini par $0 \leq x \leq 1$ et $x \leq y \leq 2x$.

\begin{parts}
    \part Dessinez le domaine d'intégration $D$ dans le plan $(x, y)$.
    \begin{solutionbox}{18cm}
        Le domaine $D$ est délimité par :
        \begin{itemize}
            \item $x = 0$ (ligne verticale)
            \item $x = 1$ (ligne verticale)
            \item $y = x$ (droite passant par l'origine)
            \item $y = 2x$ (droite passant par l'origine avec pente 2)
        \end{itemize}
        
        Le domaine est un triangle entre ces lignes pour $x \in [0, 1]$.
        Pour $x = 0$, on a $y \in [0, 0]$ (un point).
        Pour $x = 1$, on a $y \in [1, 2]$.
        Le domaine est la région entre les droites $y = x$ et $y = 2x$ pour $x \in [0, 1]$.
        
        \begin{center}
            \begin{tikzpicture}
                \begin{axis}[
                    axis lines=middle,
                    xlabel={$x$},
                    ylabel={$y$},
                    xmin=-0.2, xmax=1.2,
                    ymin=-0.2, ymax=2.2,
                    domain=0:1,
                    samples=100,
                    width=12cm,
                    height=8cm,
                    legend pos=north west
                ]
                    % Définir les chemins pour le remplissage
                    \addplot[name path=lower, domain=0:1, draw=none] {x};
                    \addplot[name path=upper, domain=0:1, draw=none] {2*x};
                    
                    % Remplir le domaine entre les deux courbes
                    \addplot[fill=gray!30, draw=none] fill between[of=lower and upper];
                    
                    % Tracer les droites
                    \addplot[domain=0:1, thick, blue] {x};
                    \addplot[domain=0:1, thick, red] {2*x};
                    
                    % Tracer les lignes verticales
                    \draw[thick, dashed] (axis cs:0,0) -- (axis cs:0,2);
                    \draw[thick, dashed] (axis cs:1,0) -- (axis cs:1,2);
                    
                    % Légende
                    \addlegendentry{$y = x$}
                    \addlegendentry{$y = 2x$}
                \end{axis}
            \end{tikzpicture}
        \end{center}
    \end{solutionbox}
    
    \part Effectuez le calcul de cette intégrale dans l'ordre $dy \, dx$.
    \begin{solutionbox}{21cm}
        \begin{align*}
            \iint_D xy \, dx \, dy &= \int_0^1 \int_x^{2x} xy \, dy \, dx
        \end{align*}
        
        Pour l'intégrale intérieure, on intègre par rapport à $y$ en traitant $x$ comme une constante :
        \begin{align*}
            \int_x^{2x} xy \, dy &= x \int_x^{2x} y \, dy\\
            &= x \left[ \frac{y^2}{2} \right]_{y=x}^{y=2x}\\
            &= x \left( \frac{(2x)^2}{2} - \frac{x^2}{2} \right)\\
            &= x \left( \frac{4x^2}{2} - \frac{x^2}{2} \right)
            = x \cdot \frac{3x^2}{2}
            = \frac{3x^3}{2}
        \end{align*}
        
        Donc :
        \begin{align*}
            \iint_D xy \, dx \, dy &= \int_0^1 \frac{3x^3}{2} \, dx\\
            &= \frac{3}{2} \int_0^1 x^3 \, dx\\
            &= \frac{3}{2} \left[ \frac{x^4}{4} \right]_0^1
            = \frac{3}{2} \cdot \frac{1}{4}
            = \frac{3}{8}
        \end{align*}
    \end{solutionbox}
    
    \part Changez l'ordre d'intégration et recalculez l'intégrale dans l'ordre $dx \, dy$.
    \begin{solutionbox}{21cm}
        Pour changer l'ordre d'intégration, il faut exprimer le domaine en fonction de $y$ d'abord.
        
        Le domaine $D$ est délimité par :
        \begin{itemize}
            \item $y = x$ (donc $x = y$)
            \item $y = 2x$ (donc $x = y/2$)
            \item $x = 0$ et $x = 1$
        \end{itemize}
        
        Pour un $y$ fixé, les bornes en $x$ sont :
        \begin{itemize}
            \item Pour $y \in [0, 1]$ : $x$ varie de $y/2$ (ligne $y = 2x$) à $y$ (ligne $y = x$)
            \item Pour $y \in [1, 2]$ : $x$ varie de $y/2$ (ligne $y = 2x$) à $1$ (ligne $x = 1$)
        \end{itemize}
        
        Donc :
        \begin{align*}
            \iint_D xy \, dx \, dy &= \int_0^1 \int_{y/2}^y xy \, dx \, dy + \int_1^2 \int_{y/2}^1 xy \, dx \, dy
        \end{align*}
        
        Pour la première intégrale ($y \in [0, 1]$) :
        \begin{align*}
            \int_{y/2}^y xy \, dx &= y \int_{y/2}^y x \, dx
            = y \left[ \frac{x^2}{2} \right]_{x=y/2}^{x=y}\\
            &= y \left( \frac{y^2}{2} - \frac{(y/2)^2}{2} \right)
            = y \left( \frac{y^2}{2} - \frac{y^2}{8} \right)
            = y \cdot \frac{3y^2}{8}
            = \frac{3y^3}{8}
        \end{align*}
        
        Pour la deuxième intégrale ($y \in [1, 2]$) :
        \begin{align*}
            \int_{y/2}^1 xy \, dx &= y \int_{y/2}^1 x \, dx
            = y \left[ \frac{x^2}{2} \right]_{x=y/2}^{x=1}\\
            &= y \left( \frac{1}{2} - \frac{(y/2)^2}{2} \right)
            = y \left( \frac{1}{2} - \frac{y^2}{8} \right)
            = \frac{y}{2} - \frac{y^3}{8}
        \end{align*}
        
        Donc :
        \begin{align*}
            \iint_D xy \, dx \, dy &= \int_0^1 \frac{3y^3}{8} \, dy + \int_1^2 \left( \frac{y}{2} - \frac{y^3}{8} \right) \, dy\\
            &= \frac{3}{8} \int_0^1 y^3 \, dy + \frac{1}{2} \int_1^2 y \, dy - \frac{1}{8} \int_1^2 y^3 \, dy\\
            &= \frac{3}{8} \left[ \frac{y^4}{4} \right]_0^1 + \frac{1}{2} \left[ \frac{y^2}{2} \right]_1^2 - \frac{1}{8} \left[ \frac{y^4}{4} \right]_1^2\\
            &= \frac{3}{8} \cdot \frac{1}{4} + \frac{1}{2} \left( \frac{4}{2} - \frac{1}{2} \right) - \frac{1}{8} \left( \frac{16}{4} - \frac{1}{4} \right)\\
            &= \frac{3}{32} + \frac{1}{2} \cdot \frac{3}{2} - \frac{1}{8} \cdot \frac{15}{4}
            = \frac{3}{8}
        \end{align*}
    \end{solutionbox}
\end{parts}
